% GENERATED FILE. Edit canonical lessons, then run npm run generate:books.
% canonical-source-hash: fnv1a64:352b4b8241d0fe9c
% canonical-lessons: KA-C115-karana, KA-C115-adudarinda, KA-C115-adakke, KA-C115-yakendare, KA-C115-mamca

\chapter{Reason, Therefore, That is why, Because, Bed}
\label{ch:ka-reason-therefore-that-is-why-because-bed}
\hlchaptermodality{115}

\begin{chapteropening}
\textbf{By the end of this chapter:} \emph{I can say a reason, therefore, that is why, because and a bed.}

Complete the last lesson of chapter 115: I can say a reason, therefore, that is why, because and a bed.
\end{chapteropening}

\section[\texorpdfstring{\kn{ಕಾರಣ} (kāraṇa)}{ಕಾರಣ (kāraṇa)}]{\kn{ಕಾರಣ} (kāraṇa) — a reason}

\label{lesson:KA-C115-karana}



\begin{quote}
Before the new one: say the Kannada for a decision, then the Kannada for to wish for.
\end{quote}

\subsection*{You'll want to know: \kn{ಕಾರಣ}}
\textbf{\kn{ಕಾರಣ}} — \emph{kāraṇa} — \textquotedblleft{}a reason\textquotedblright{}.

\begin{cousinweb}
\textbf{\kn{ಕಾರಣ}} (\emph{kāraṇa}) is Sanskrit \emph{kāraṇa}, \textquotedblleft{}cause\textquotedblright{}, taken over whole as a learned loan, a tatsama (CDIAL 3057).
\end{cousinweb}

\begin{grammarlens}[title={What you've built}]
One more word for this chapter.
\end{grammarlens}

\subsection*{Your turn}
\begin{itemize}
\raggedright
  \item \emph{Say it:} \emph{kāraṇa}
  \item \emph{Say it:} \emph{kāraṇa}, once more
  \item \emph{Say it:} say \emph{kāraṇa}
  \item \emph{Recall:} say the Kannada for a decision, then the Kannada for to wish for, then say \emph{kāraṇa} again
\end{itemize}

\begin{tcolorbox}[breakable,colback=teal!4,colframe=teal!35!black,title={Before you move on}]
What does \kn{ಕಾರಣ} mean? (\textquotedblleft{}A reason\textquotedblright{}.) Say it once more.
\end{tcolorbox}
\section[\texorpdfstring{\kn{ಆದುದರಿಂದ} (ādudarinda)}{ಆದುದರಿಂದ (ādudarinda)}]{\kn{ಆದುದರಿಂದ} (ādudarinda) — therefore}

\label{lesson:KA-C115-adudarinda}



\begin{quote}
Before the new one: say the Kannada for to wish for, then the Kannada for a reason.
\end{quote}

\subsection*{You'll want to know: \kn{ಆದುದರಿಂದ}}
\textbf{\kn{ಆದುದರಿಂದ}} — \emph{ādudarinda} — \textquotedblleft{}therefore\textquotedblright{}.

\begin{cousinweb}
\textbf{\kn{ಆದುದರಿಂದ}} (\emph{ādudarinda}) is native Dravidian: \emph{ādudu}, \textquotedblleft{}that which became\textquotedblright{}, from \emph{āgu}, \textquotedblleft{}to become\textquotedblright{} (DEDR 333), with \emph{adu}, \textquotedblleft{}that\textquotedblright{} (DEDR 1), and \emph{-inda}, \textquotedblleft{}from, because of\textquotedblright{}: \textquotedblleft{}because of that which happened\textquotedblright{}. DEDR 333 compares Sanskrit \emph{ām}.
\end{cousinweb}

\begin{grammarlens}[title={What you've built}]
One more word for this chapter.
\end{grammarlens}

\subsection*{Your turn}
\begin{itemize}
\raggedright
  \item \emph{Say it:} \emph{ādudarinda}
  \item \emph{Say it:} \emph{ādudarinda}, once more
  \item \emph{Say it:} say \emph{ādudarinda}
  \item \emph{Recall:} say the Kannada for to wish for, then the Kannada for a reason, then say \emph{ādudarinda} again
\end{itemize}

\begin{tcolorbox}[breakable,colback=teal!4,colframe=teal!35!black,title={Before you move on}]
What does \kn{ಆದುದರಿಂದ} mean? (\textquotedblleft{}Therefore\textquotedblright{}.) Say it once more.
\end{tcolorbox}
\section[\texorpdfstring{\kn{ಅದಕ್ಕೇ} (adakkē)}{ಅದಕ್ಕೇ (adakkē)}]{\kn{ಅದಕ್ಕೇ} (adakkē) — that is why}

\label{lesson:KA-C115-adakke}



\begin{quote}
Before the new one: say the Kannada for a reason, then the Kannada for therefore.
\end{quote}

\subsection*{You'll want to know: \kn{ಅದಕ್ಕೇ}}
\textbf{\kn{ಅದಕ್ಕೇ}} — \emph{adakkē} — \textquotedblleft{}that is why\textquotedblright{}.

\begin{cousinweb}
\textbf{\kn{ಅದಕ್ಕೇ}} (\emph{adakkē}) is native Dravidian: \emph{adu}, \textquotedblleft{}that\textquotedblright{} (DEDR 1), with \emph{-kke}, \textquotedblleft{}for\textquotedblright{}, and \emph{-ē}, \textquotedblleft{}exactly\textquotedblright{}: \textquotedblleft{}exactly for that\textquotedblright{}.
\end{cousinweb}

\begin{grammarlens}[title={What you've built}]
One more word for this chapter.
\end{grammarlens}

\subsection*{Your turn}
\begin{itemize}
\raggedright
  \item \emph{Say it:} \emph{adakkē}
  \item \emph{Say it:} \emph{adakkē}, once more
  \item \emph{Say it:} say \emph{adakkē}
  \item \emph{Recall:} say the Kannada for a reason, then the Kannada for therefore, then say \emph{adakkē} again
\end{itemize}

\begin{tcolorbox}[breakable,colback=teal!4,colframe=teal!35!black,title={Before you move on}]
What does \kn{ಅದಕ್ಕೇ} mean? (\textquotedblleft{}That is why\textquotedblright{}.) Say it once more.
\end{tcolorbox}
\section[\texorpdfstring{\kn{ಯಾಕೆಂದರೆ} (yākendare)}{ಯಾಕೆಂದರೆ (yākendare)}]{\kn{ಯಾಕೆಂದರೆ} (yākendare) — because}

\label{lesson:KA-C115-yakendare}



\begin{quote}
Before the new one: say the Kannada for therefore, then the Kannada for that is why.
\end{quote}

\subsection*{You'll want to know: \kn{ಯಾಕೆಂದರೆ}}
\textbf{\kn{ಯಾಕೆಂದರೆ}} — \emph{yākendare} — \textquotedblleft{}because\textquotedblright{}.

\begin{cousinweb}
\textbf{\kn{ಯಾಕೆಂದರೆ}} (\emph{yākendare}) is native Dravidian: \emph{yāke}, \textquotedblleft{}why\textquotedblright{} (DEDR 5151), and \emph{endare}, \textquotedblleft{}if one says\textquotedblright{}, from \emph{en-}, \textquotedblleft{}to say\textquotedblright{} (DEDR 868). \textquotedblleft{}If you say why\textquotedblright{} introduces the reason.
\end{cousinweb}

\begin{grammarlens}[title={What you've built}]
One more word for this chapter.
\end{grammarlens}

\subsection*{Your turn}
\begin{itemize}
\raggedright
  \item \emph{Say it:} \emph{yākendare}
  \item \emph{Say it:} \emph{yākendare}, once more
  \item \emph{Say it:} say \emph{yākendare}
  \item \emph{Recall:} say the Kannada for therefore, then the Kannada for that is why, then say \emph{yākendare} again
\end{itemize}

\begin{tcolorbox}[breakable,colback=teal!4,colframe=teal!35!black,title={Before you move on}]
What does \kn{ಯಾಕೆಂದರೆ} mean? (\textquotedblleft{}Because\textquotedblright{}.) Say it once more.
\end{tcolorbox}
\section[\texorpdfstring{\kn{ಮಂಚ} (maṁca)}{ಮಂಚ (maṁca)}]{\kn{ಮಂಚ} (maṁca) — a bed}

\label{lesson:KA-C115-mamca}



\begin{quote}
Before the new one: say the Kannada for that is why, then the Kannada for because.
\end{quote}

\subsection*{You'll want to know: \kn{ಮಂಚ}}
\textbf{\kn{ಮಂಚ}} — \emph{maṁca} — \textquotedblleft{}a bed\textquotedblright{}.

\begin{cousinweb}
\textbf{\kn{ಮಂಚ}} (\emph{maṁca}) is Sanskrit \emph{mañca}, \textquotedblleft{}stage, platform; bed\textquotedblright{}, taken over whole as a learned loan, a tatsama (CDIAL 9715).
\end{cousinweb}

\begin{grammarlens}[title={What you've built}]
One more word for this chapter.
\end{grammarlens}

\subsection*{Your turn}
\begin{itemize}
\raggedright
  \item \emph{Say it:} \emph{maṁca}
  \item \emph{Say it:} \emph{maṁca}, once more
  \item \emph{Say it:} say \emph{maṁca}
  \item \emph{Recall:} say the Kannada for that is why, then the Kannada for because, then say \emph{maṁca} again
\end{itemize}

\begin{tcolorbox}[breakable,colback=teal!4,colframe=teal!35!black,title={Before you move on}]
What does \kn{ಮಂಚ} mean? (\textquotedblleft{}A bed\textquotedblright{}.) Say it once more.
\end{tcolorbox}
